% GENERATED FILE. Edit canonical lessons, then run npm run generate:books.
% canonical-source-hash: fnv1a64:40f5270d53b1f2db
% canonical-lessons: LA-C203-dubito, LA-C203-quaero, LA-C203-consumo, LA-C203-vexo, LA-C203-obliviscor

\chapter{To doubt, To look for, To use up, To annoy, To forget}
\label{ch:la-to-doubt-to-look-for-to-use-up-to-annoy-to-forget}
\hlchaptermodality{203}

\begin{chapteropening}
\textbf{By the end of this chapter:} \emph{I can say to doubt, to look for, to use up, to annoy and to forget.}

Complete the last lesson of chapter 203: I can say to doubt, to look for, to use up, to annoy and to forget.
\end{chapteropening}

\section[dubitō, dubitāre]{dubitō, dubitāre — to doubt, to hesitate}

\label{lesson:LA-C203-dubito}



\begin{quote}
Before the new one: say the Latin for to do, then the Latin for to go forward.
\end{quote}

\subsection*{You'll want to know: dubitō, dubitāre}
\textbf{dubitō, dubitāre} — \textquotedblleft{}to doubt, to hesitate\textquotedblright{}.

\begin{cousinweb}
\textbf{dubitō} is built on \emph{dubius}, \textquotedblleft{}doubtful\textquotedblright{} (Lewis \& Short), whose first part is \emph{duo}, \textquotedblleft{}two\textquotedblright{}: to be of two minds.
\end{cousinweb}

\begin{grammarlens}[title={What you've built}]
One more word for this chapter.
\end{grammarlens}

\subsection*{Your turn}
\begin{itemize}
\raggedright
  \item \emph{Say it:} \emph{dubitō, dubitāre}
  \item \emph{Say it:} \emph{dubitō, dubitāre}, once more
  \item \emph{Say it:} say \emph{dubitō, dubitāre}
  \item \emph{Recall:} say the Latin for to do, then the Latin for to go forward, then say \emph{dubitō, dubitāre} again
\end{itemize}

\begin{tcolorbox}[breakable,colback=teal!4,colframe=teal!35!black,title={Before you move on}]
What does dubitō, dubitāre mean? (\textquotedblleft{}To doubt, to hesitate\textquotedblright{}.) Say it once more.
\end{tcolorbox}
\section[quaerō, quaerere]{quaerō, quaerere — to look for, to ask}

\label{lesson:LA-C203-quaero}



\begin{quote}
Before the new one: say the Latin for to go forward, then the Latin for to doubt.
\end{quote}

\subsection*{You'll want to know: quaerō, quaerere}
\textbf{quaerō, quaerere} — \textquotedblleft{}to look for, to ask\textquotedblright{}.

\begin{cousinweb}
\textbf{quaerō} has no sure origin. Old Latin wrote it \emph{quairō}, and \emph{quaesō}, \textquotedblleft{}I beg\textquotedblright{}, is an older form of the same verb (Lewis \& Short).
\end{cousinweb}

\begin{grammarlens}[title={What you've built}]
One more word for this chapter.
\end{grammarlens}

\subsection*{Your turn}
\begin{itemize}
\raggedright
  \item \emph{Say it:} \emph{quaerō, quaerere}
  \item \emph{Say it:} \emph{quaerō, quaerere}, once more
  \item \emph{Say it:} say \emph{quaerō, quaerere}
  \item \emph{Recall:} say the Latin for to go forward, then the Latin for to doubt, then say \emph{quaerō, quaerere} again
\end{itemize}

\begin{tcolorbox}[breakable,colback=teal!4,colframe=teal!35!black,title={Before you move on}]
What does quaerō, quaerere mean? (\textquotedblleft{}To look for, to ask\textquotedblright{}.) Say it once more.
\end{tcolorbox}
\section[cōnsūmō, cōnsūmere]{cōnsūmō, cōnsūmere — to use up}

\label{lesson:LA-C203-consumo}



\begin{quote}
Before the new one: say the Latin for to doubt, then the Latin for to look for.
\end{quote}

\subsection*{You'll want to know: cōnsūmō, cōnsūmere}
\textbf{cōnsūmō, cōnsūmere} — \textquotedblleft{}to use up\textquotedblright{}.

\begin{cousinweb}
\textbf{cōnsūmō} is \emph{con-}, \textquotedblleft{}completely\textquotedblright{}, with \emph{sūmō}, \textquotedblleft{}to take\textquotedblright{} (EtymDB 2.1): to take wholly, to use up.
\end{cousinweb}

\begin{grammarlens}[title={What you've built}]
One more word for this chapter.
\end{grammarlens}

\subsection*{Your turn}
\begin{itemize}
\raggedright
  \item \emph{Say it:} \emph{cōnsūmō, cōnsūmere}
  \item \emph{Say it:} \emph{cōnsūmō, cōnsūmere}, once more
  \item \emph{Say it:} say \emph{cōnsūmō, cōnsūmere}
  \item \emph{Recall:} say the Latin for to doubt, then the Latin for to look for, then say \emph{cōnsūmō, cōnsūmere} again
\end{itemize}

\begin{tcolorbox}[breakable,colback=teal!4,colframe=teal!35!black,title={Before you move on}]
What does cōnsūmō, cōnsūmere mean? (\textquotedblleft{}To use up\textquotedblright{}.) Say it once more.
\end{tcolorbox}
\section[vexō, vexāre]{vexō, vexāre — to annoy, to harass}

\label{lesson:LA-C203-vexo}



\begin{quote}
Before the new one: say the Latin for to look for, then the Latin for to use up.
\end{quote}

\subsection*{You'll want to know: vexō, vexāre}
\textbf{vexō, vexāre} — \textquotedblleft{}to annoy, to harass\textquotedblright{}.

\begin{cousinweb}
\textbf{vexō} is usually taken as a frequentative of \emph{vehō}, \textquotedblleft{}to carry\textquotedblright{} (Lewis \& Short): to keep carrying about, to toss, so to harass.
\end{cousinweb}

\begin{grammarlens}[title={What you've built}]
One more word for this chapter.
\end{grammarlens}

\subsection*{Your turn}
\begin{itemize}
\raggedright
  \item \emph{Say it:} \emph{vexō, vexāre}
  \item \emph{Say it:} \emph{vexō, vexāre}, once more
  \item \emph{Say it:} say \emph{vexō, vexāre}
  \item \emph{Recall:} say the Latin for to look for, then the Latin for to use up, then say \emph{vexō, vexāre} again
\end{itemize}

\begin{tcolorbox}[breakable,colback=teal!4,colframe=teal!35!black,title={Before you move on}]
What does vexō, vexāre mean? (\textquotedblleft{}To annoy, to harass\textquotedblright{}.) Say it once more.
\end{tcolorbox}
\section[oblīvīscor, oblīvīscī]{oblīvīscor, oblīvīscī — to forget}

\label{lesson:LA-C203-obliviscor}



\begin{quote}
Before the new one: say the Latin for to use up, then the Latin for to annoy.
\end{quote}

\subsection*{You'll want to know: oblīvīscor, oblīvīscī}
\textbf{oblīvīscor, oblīvīscī} — \textquotedblleft{}to forget\textquotedblright{}.

\begin{cousinweb}
\textbf{oblīvīscor} is usually explained as \emph{ob}, \textquotedblleft{}over\textquotedblright{}, with a root for smearing or smoothing: the mind wiped over. Lewis \& Short tie it instead to \emph{līvor}, \textquotedblleft{}a dark colour\textquotedblright{}.
\end{cousinweb}

\begin{grammarlens}[title={What you've built}]
One more word for this chapter.
\end{grammarlens}

\subsection*{Your turn}
\begin{itemize}
\raggedright
  \item \emph{Say it:} \emph{oblīvīscor, oblīvīscī}
  \item \emph{Say it:} \emph{oblīvīscor, oblīvīscī}, once more
  \item \emph{Say it:} say \emph{oblīvīscor, oblīvīscī}
  \item \emph{Recall:} say the Latin for to use up, then the Latin for to annoy, then say \emph{oblīvīscor, oblīvīscī} again
\end{itemize}

\begin{tcolorbox}[breakable,colback=teal!4,colframe=teal!35!black,title={Before you move on}]
What does oblīvīscor, oblīvīscī mean? (\textquotedblleft{}To forget\textquotedblright{}.) Say it once more.
\end{tcolorbox}
